% ============================================================
%  Technical Note: Series Inductance Design in DAB and CLLC
%  DC-DC Isolated Converters
%  Author: Davide
% ============================================================

\documentclass[a4paper, 11pt]{article}

\usepackage[T1]{fontenc}
\usepackage[utf8]{inputenc}
\usepackage[english]{babel}
\usepackage{amsmath, amssymb}
\usepackage{booktabs}
\usepackage{geometry}
\usepackage[dvipsnames]{xcolor}
\usepackage{hyperref}
\usepackage{lmodern}
\usepackage{microtype}
\usepackage{parskip}

\geometry{margin=2.5cm}

\hypersetup{
	colorlinks = true,
	linkcolor  = blue!60!black,
	urlcolor   = blue!60!black,
}

\title{\textbf{Series Inductance Design in Isolated DC-DC Converters:\\
		DAB vs.\ CLLC}}
\author{Davide}
\date{\today}

% ============================================================
\begin{document}
	\maketitle
	\tableofcontents
	\bigskip
	
	% ------------------------------------------------------------
	\section{Introduction}
	
	Both the Dual Active Bridge (DAB) and the CLLC resonant converter are
	bidirectional, galvanically isolated DC-DC topologies widely used in
	energy-storage interfaces, on-board chargers, and solid-state
	transformer applications.  In both cases a series inductance
	$L_s$—realised either as the leakage inductance of the transformer or
	as an external component—plays a central role in limiting and
	controlling the power flow.  The underlying physics, however, differs
	substantially between the two topologies, leading to different design
	equations and different degrees of freedom for the power-stage
	engineer.
	
	This note derives the sizing equation for $L_s$ in both converters,
	assuming a unity transformer turns-ratio ($n = 1$) throughout, and
	draws a direct comparison between the two approaches.
	
	% ------------------------------------------------------------
	\section{Dual Active Bridge (DAB)}
	
	\subsection{Operating Principle}
	
	The DAB consists of two full-bridge voltage-source inverters connected
	through a high-frequency transformer and its series inductance $L_s$.
	The two bridges generate square-wave voltages at the same switching
	frequency $f_s$, and power is transferred by introducing a controlled
	phase-shift $\phi$ (in radians) between the two bridge outputs.  With
	Single-Phase-Shift (SPS) modulation and $n = 1$, the average
	transferred power is~\cite{dedoncker1991}:
	
	\begin{equation}
		P = \frac{V_1 \, V_2}{2\pi f_s L_s}
		\cdot \phi \!\left(1 - \frac{|\phi|}{\pi}\right),
		\qquad |\phi| \leq \frac{\pi}{2}.
		\label{eq:dab_power}
	\end{equation}
	
	\subsection{Maximum Power and Inductance Sizing}
	
	Differentiating \eqref{eq:dab_power} with respect to $\phi$ shows that
	the power is maximised at $\phi = \pi/2$, giving:
	
	\begin{equation}
		P_{\max} = \frac{V_1 \, V_2}{8\, f_s\, L_s}.
		\label{eq:dab_pmax}
	\end{equation}
	
	Inverting \eqref{eq:dab_pmax} yields the fundamental sizing equation:
	
	\begin{equation}
		\boxed{L_s = \frac{V_1 \, V_2}{8\, f_s\, P_{\max}}.}
		\label{eq:dab_ls}
	\end{equation}
	
	\subsection{Design Considerations}
	
	Equation~\eqref{eq:dab_ls} is \emph{closed-form}: given the DC bus
	voltages, the switching frequency, and the rated power, $L_s$ is
	determined uniquely.  Larger $L_s$ reduces $P_{\max}$ and increases
	the phase-shift required at a given operating point, which can
	improve light-load ZVS margins but worsens reactive-current stress at
	full load.  The relationship between power and phase-shift is
	\emph{linear} in $\phi$ for small angles, and the converter behaves
	essentially as a controlled current source driven by a voltage
	difference.
	
	% ------------------------------------------------------------
	\section{CLLC Resonant Converter}
	
	\subsection{Operating Principle}
	
	The CLLC is the bidirectional extension of the LLC resonant converter.
	The resonant network is formed by the series inductance $L_r$,
	the series capacitance $C_r$, and the magnetising inductance $L_m$ on
	the primary side; a symmetric network ($L_r'$, $C_r'$) is placed on
	the secondary to preserve bidirectional resonant behaviour.
	The converter is operated close to the series resonant frequency:
	
	\begin{equation}
		f_r = \frac{1}{2\pi\sqrt{L_r C_r}}.
		\label{eq:cllc_fr}
	\end{equation}
	
	\subsection{First Harmonic Approximation (FHA)}
	
	Under FHA, the nonlinear switching network is replaced by sinusoidal
	sources and the load is represented by its equivalent AC resistance.
	For a resistive output load $R_{\text{out}}$ and $n = 1$, the
	equivalent AC load resistance seen by the resonant tank is:
	
	\begin{equation}
		R_{ac} = \frac{8}{\pi^2} \cdot \frac{V_2^2}{P}
		= \frac{8}{\pi^2} \cdot R_{\text{out}}.
		\label{eq:cllc_rac}
	\end{equation}
	
	\subsection{Quality Factor and Inductance Sizing}
	
	The characteristic impedance of the resonant tank is:
	
	\begin{equation}
		Z_r = \sqrt{\frac{L_r}{C_r}},
		\label{eq:cllc_zr}
	\end{equation}
	
	and the loaded quality factor at resonance is:
	
	\begin{equation}
		Q = \frac{Z_r}{R_{ac}} = \frac{\sqrt{L_r / C_r}}{R_{ac}}.
		\label{eq:cllc_Q}
	\end{equation}
	
	From \eqref{eq:cllc_fr} and \eqref{eq:cllc_Q}, solving for $L_r$ and
	$C_r$:
	
	\begin{equation}
		\boxed{L_r = \frac{Q \cdot R_{ac}}{2\pi f_r},}
		\qquad
		C_r = \frac{1}{Q \cdot R_{ac} \cdot 2\pi f_r}.
		\label{eq:cllc_LrCr}
	\end{equation}
	
	A typical design targets $Q \leq 0.5 \div 1$ at full load in order
	to guarantee ZVS over an adequate load range.
	
	\subsection{Design Considerations}
	
	Unlike the DAB, the CLLC does not have a single closed-form equation
	that directly links $L_r$ to $P_{\max}$.  The role of $L_r$ is to set
	the tank impedance $Z_r$ and, together with $L_m$, to shape the
	converter voltage gain characteristic.  The three key parameters—$f_r$,
	$Q$, and the inductance ratio $k = L_m / L_r$—must be chosen
	simultaneously to satisfy: (i) the required voltage-gain range;
	(ii) ZVS over the desired load range; and (iii) an acceptable
	peak-current stress.  The primary control variable is the switching
	frequency $f_s$, which is varied around $f_r$ to regulate the output
	voltage.
	
	% ------------------------------------------------------------
	\section{Direct Comparison}
	
	Table~\ref{tab:comparison} summarises the principal differences between
	the two approaches.
	
	\begin{table}[ht]
		\centering
		\caption{DAB vs.\ CLLC: comparison of the series inductance design.}
		\label{tab:comparison}
		\renewcommand{\arraystretch}{1.4}
		\begin{tabular}{lll}
			\toprule
			\textbf{Aspect} & \textbf{DAB} & \textbf{CLLC} \\
			\midrule
			Power-transfer mechanism
			& Phase-shift across $L_s$
			& Series resonance $L_r$–$C_r$ \\
			$P_{\max}$ sizing equation
			& $L_s = \dfrac{V_1 V_2}{8 f_s P_{\max}}$
			& $L_r = \dfrac{Q \cdot R_{ac}}{2\pi f_r}$ \\
			Primary control variable
			& Phase-shift $\phi$ at fixed $f_s$
			& Frequency ratio $f_s / f_r$ \\
			ZVS condition
			& Conditional (depends on $\phi$, load)
			& Natural around $f_r$ \\
			Effect of larger $L_s$
			& Lower $P_{\max}$, higher $\phi$ at rated load
			& Higher $Z_r$, higher $Q$ \\
			Design complexity
			& Closed-form, single equation
			& Iterative, multi-parameter \\
			\bottomrule
		\end{tabular}
	\end{table}
	
	% ------------------------------------------------------------
	\section{Summary}
	
	For the DAB, the series inductance $L_s$ is determined by a single,
	closed-form expression \eqref{eq:dab_ls} that directly links the
	maximum transferable power to the DC bus voltages and the switching
	frequency.  The design procedure is therefore straightforward.
	
	For the CLLC, the resonant inductance $L_r$ is obtained from
	\eqref{eq:cllc_LrCr} once the resonant frequency $f_r$ and the target
	quality factor $Q$ have been chosen.  The choice of $Q$ and of the
	inductance ratio $k = L_m/L_r$ is, however, inherently iterative
	because these parameters jointly determine the voltage-gain curve, the
	ZVS range, and the peak circulating current.
	
	% ------------------------------------------------------------
	\begin{thebibliography}{9}
		\bibitem{dedoncker1991}
		R.~W.~De~Doncker, D.~M.~Divan, and M.~H.~Kheraluwala,
		``A three-phase soft-switched high-power-density DC/DC converter
		for high-power applications,''
		\emph{IEEE Transactions on Industry Applications},
		vol.~27, no.~1, pp.~63--73, Jan./Feb.\ 1991.
		
		\bibitem{beiranvand2011}
		R.~Beiranvand, B.~Rashidian, M.~R.~Zolghadri, and
		S.~M.~H.~Alavi,
		``Using LLC resonant converter for designing wide-range voltage
		source,''
		\emph{IEEE Transactions on Industrial Electronics},
		vol.~58, no.~5, pp.~1746--1756, May 2011.
	\end{thebibliography}
	
\end{document}